\chapter{Pesin's entropy formula (volume case)}
\label{chap:pesin}

The Margulis--Ruelle inequality of Chapter~\ref{chap:ksentropy} bounds the Kolmogorov--Sinai
entropy of a smooth ergodic map from above by its positive Lyapunov exponents,
\(h_\mu(T)\le\sum_{\lambda_i>0}\lambda_i\). Pesin's theorem asserts that for the geometrically
distinguished \emph{SRB} (Sinai--Ruelle--Bowen) measures this bound is an \emph{equality},
\[
  h_\mu(T) \;=\; \sum_{\lambda_i>0}\lambda_i ,
\]
identifying metric entropy with the total exponential stretching along unstable directions. This
chapter formalizes the reverse inequality \(\sum_{\lambda_i>0}\lambda_i\le h_\mu(T)\) and the
resulting equality in the \emph{volume case} --- the honest fragment of the SRB condition that is
expressible without the Pesin unstable-manifold / foliation machinery Mathlib lacks --- and closes
with the first non-vacuous full-system witness, the doubling map on the circle
(\(h_\mu(T)=\sum\lambda^{+}=\log 2\)). This is the second hard leaf of the smooth theory (issue
\#10); the general mixed-spectrum Ledrappier--Young reverse inequality remains a disclosed
Mathlib-scale wall.

The upper bound is the sharp Margulis--Ruelle inequality \Cref{margulisRuelle_sharp}, which needs
\emph{no} SRB hypothesis and is developed in the classical-entropy chapter; here we recall it only
as the \(\le\) direction and concentrate on the reverse. The reverse inequality does not go through
the general Ledrappier--Strelcyn--Young route (absolute continuity of the conditional measures on
unstable manifolds, unavailable in Mathlib); instead it is discharged, in the volume case, from
\emph{Rokhlin's inequality} --- an unconditional, generator-free lower bound
\(\int\log\abs{\det D_xT}\,d\mu\le h_\mu(T)\) --- together with the trace--determinant identity and
a nonnegative-spectrum gate. Throughout, the phase space is
\(E=\operatorname{EuclideanSpace}\,\R\,(\Fin d)\), \(\mu\) is a probability measure, \(T\) is
differentiable with everywhere-nonsingular derivative, and the derivative cocycle
\(x\mapsto D_xT\) (\Cref{derivativeCocycle}) carries log-integrable data
(\Cref{IntegrableLogNorm} for \(DT\) and \(DT^{-1}\)), so that the multiplicative ergodic theorem
supplies the Lyapunov exponents \(\lambda_i\) (\Cref{exponents}).

\section{Rokhlin's inequality: the generator-free lower bound}

Rokhlin's volume-distortion \emph{identity} \(h_\mu(T,\xi)=\int\log\abs{\det D_xT}\,d\mu\)
(\Cref{ksEntropyPartition_eq_integral_log_abs_det}) needs a one-sided \emph{generating}
cell-injective partition. Dropping the generating hypothesis while keeping the injectivity
change-of-variables computation (\Cref{condEntropy_comap_eq_integral_log_abs_det}) turns the
identity into a one-sided, partition-free \emph{inequality}, which is exactly the reverse half of
Pesin's theorem in the absolutely continuous case. The key structural fact making this
unconditional is that the strict-future \(\sigma\)-algebra of \emph{any} partition sits below the
comap of \(T\).

\begin{lemma}[Strict future lies below \(\operatorname{comap}T\)]\label{strictFuture_le_comap}
  \lean{ErgodicTheory.strictFuture_le_comap}\leanok
  \uses{ksJoin}
  For a measure-preserving \(T\) and a finite measurable partition \(\xi\), the strict-future
  \(\sigma\)-algebra generated by the pullbacks of the iterated joins lies below the comap of
  \(T\):
  \[
    \bigvee_{k}\sigma\bigl((\textstyle\bigvee_{j<k}T^{-j}\xi)\text{-pullback}\bigr)
      \;\le\; \operatorname{comap}T\,m_E .
  \]
\end{lemma}
\begin{proof}\leanok
  \uses{ksJoin}
  This is the generator-free half of the equality
  \(\bigvee_k\sigma(\dots)=\operatorname{comap}T\,m_E\) that holds for a one-sided
  \emph{generating} partition: the reverse containment needs generation, but \(\le\) is
  unconditional. Per term the pulled-back generated \(\sigma\)-algebra equals the comap of \(T\)
  applied to the join's own generated \(\sigma\)-algebra, which is below \(m_E\); \texttt{iSup\_le}
  and monotonicity of comap close the bound. This \(\le\) is precisely what removes the generator
  hypothesis from Rokhlin's inequality below.
\end{proof}

\begin{theorem}[Rokhlin's inequality]\label{integral_log_abs_det_le_ksEntropy}
  \lean{ErgodicTheory.integral_log_abs_det_le_ksEntropy}\leanok
  \uses{strictFuture_le_comap, ksEntropy, ksEntropyPartition, IsInjectivityPartition, condEntropy_comap_eq_integral_log_abs_det, ksEntropyPartition_eq_condEntropy_iSup}
  Let \(T\) be a measure-preserving, differentiable self-map of \(\R^d\) with
  \(\mu\ll\operatorname{volume}\) and everywhere-nonsingular derivative, let \(\xi\) be a nonempty
  injectivity partition (\Cref{IsInjectivityPartition}), and assume \(\log\rho\in L^1(\mu)\) for the
  density \(\rho=d\mu/d\operatorname{vol}\) and \(\log\abs{\det DT}\in L^1(\mu)\). Then, with
  \emph{no} generating-partition hypothesis,
  \[
    \int \log\,\abs{\det D_xT}\;d\mu(x) \;\le\; h_\mu(T) .
  \]
  This is the generator-free lower bound on entropy (Rokhlin 1967, \S9; Parry 1969; Coud\`ene,
  \emph{Ergodic Theory and Dynamical Systems}, Cor.~12.1): the integrated log volume distortion
  never exceeds the Kolmogorov--Sinai entropy. It needs no ergodicity and no expansion --- those
  enter only for Rokhlin's \emph{equality}. Here \(\mu\ll\operatorname{volume}\) enters solely
  through Coud\`ene's Remark \(\abs{T'_\mu}=\abs{\det DT}\) a.e., identifying the Radon--Nikodym
  distortion with the Jacobian determinant. The hypotheses are non-vacuous on \(\R^d\): e.g.
  \(T=\operatorname{id}\) with any a.c. probability \(\mu\) gives \(0\le h_\mu(\operatorname{id})\).
\end{theorem}
\begin{proof}\leanok
  \uses{strictFuture_le_comap, condEntropy_comap_eq_integral_log_abs_det, ksEntropyPartition_eq_condEntropy_iSup, ksEntropy}
  Four steps chain. First, the injectivity change-of-variables computation
  (\Cref{condEntropy_comap_eq_integral_log_abs_det}) evaluates the conditional entropy of \(\xi\)
  against \(\operatorname{comap}T\,m_E\) to \(\int\log\abs{\det D_xT}\,d\mu\), \emph{independently of
  the partition}. Second, conditioning is antitone in the \(\sigma\)-algebra, so
  \Cref{strictFuture_le_comap} bounds \(H(\xi\mid\operatorname{comap}T\,m_E)\) above by
  \(H(\xi\mid\mathcal S_\infty)\) against the larger strict-future \(\sigma\)-algebra
  \(\mathcal S_\infty\). Third, the sharp-rate identity
  (\Cref{ksEntropyPartition_eq_condEntropy_iSup}) identifies \(H(\xi\mid\mathcal S_\infty)\) with
  the per-partition entropy \(h_\mu(T,\xi)\). Finally \(h_\mu(T,\xi)\le h_\mu(T)\) lifts to the
  system entropy, and the fixed left-hand side survives all three monotone steps.
\end{proof}

\section{The SRB property and the unstable-Jacobian bridge}

Pesin's formula holds \emph{exactly} for SRB measures: those whose conditional measures on unstable
manifolds are absolutely continuous with respect to leaf (Riemannian) volume
(Ledrappier--Strelcyn--Young). In full generality that condition is not stateable in Mathlib --- it
needs the Pesin unstable-manifold theorem integrating the measurable distribution \(E^u\) into the
foliation \(W^u\), the disintegration of \(\mu\) along \emph{that specific} foliation, and the
leaf-volume measure, none of which exist in the library. What \emph{is} expressible is the
\emph{volume case}: the situation in which the whole space is a single unstable leaf, so leaf volume
is ambient volume and the SRB condition collapses to \(\mu\ll\operatorname{volume}\).

\begin{definition}[The SRB property, volume case]\label{SRBProperty}
  \lean{ErgodicTheory.SRBProperty}\leanok
  A map \(T\) preserves an \emph{SRB (volume-case) measure} \(\mu\) if \(\mu\) is absolutely
  continuous with respect to ambient Lebesgue volume, \(\mu\ll\operatorname{volume}\). This is the
  volume-case specialization of the Ledrappier--Strelcyn--Young SRB condition. It is a genuine,
  non-vacuous, dischargeable hypothesis: \(\operatorname{volume}\ll\operatorname{volume}\) witnesses
  inhabitation, while a Dirac mass is \emph{not} \(\ll\operatorname{volume}\). Three earlier
  framings on issue \#10 were all defective and are recorded so the rescope is auditable: an
  \(\texttt{acConditionalUnstable}:\texttt{True}\) field (vacuous, dischargeable by
  \(\langle\texttt{trivial}\rangle\)); an existential
  \(\exists\pi\,\eta,\ \operatorname{singularPart}(\operatorname{condDistrib}\dots)=0\) (also
  vacuous --- satisfied by \emph{every} \(\mu\), since an existential over an arbitrary measurable
  factor self-trivializes); and an \(\texttt{opaque}\) marker (honest but \emph{unconsumable}, so
  the reverse inequality stayed a documented \(\texttt{BLOCKED}\) leaf). The field
  \(\mu\ll\operatorname{volume}\) is the honest fix that is both genuine and \emph{consumable} ---
  the reverse inequality is proved from it.
\end{definition}

\begin{definition}[Unstable-Jacobian rate]\label{UnstableJacobianRate}
  \lean{ErgodicTheory.UnstableJacobianRate}\leanok
  \uses{exponents}
  For an ergodic differentiable \(T\) with nonsingular log-integrable derivative cocycle, the
  predicate \(\operatorname{UnstableJacobianRate}\,\chi\) asserts that a candidate integrand
  \(\chi:E\to\R\) is \(\mu\)-a.e.\ equal to the deterministic positive-exponent sum
  \(\sum_i\lambda_i^{+}=\operatorname{sumPosExp}\). Geometrically \(\chi(x)\) is meant to be the
  a.e.\ orbit growth rate of the unstable Jacobian \(\log\abs{\det (D_xT^{[n]})|_{E^u(x)}}\); since
  the spectrum is a.e.\ constant by ergodicity, that rate is a.e.\ the constant
  \(\operatorname{sumPosExp}\). This is the bridge object that phrases Pesin's formula in its
  genuine integral form \(h_\mu(T)=\int\chi\,d\mu\): a probability measure integrates the
  a.e.-constant \(\chi\) to itself, \(\int\chi\,d\mu=\operatorname{sumPosExp}\). In the volume case
  the reverse inequality is proved \emph{without} this interface (it goes through Rokhlin's
  inequality directly); the object is retained only to state the literal integral form.
\end{definition}

\section{The reverse inequality and the spectrum gate}

The reverse inequality chains Rokhlin's inequality with the trace--determinant identity, gated by a
nonnegative Lyapunov spectrum. The gate is exactly where the volume case is delimited.

\begin{lemma}[Nonnegative spectrum collapses the positive-part sum]\label{sumPosExp_eq_sumAllExp_of_nonneg}
  \lean{ErgodicTheory.sumPosExp_eq_sumAllExp_of_nonneg}\leanok
  \uses{exponents}
  If every Lyapunov exponent of the derivative cocycle is nonnegative,
  \(0\le\lambda_i\) for all \(i\), then the positive-part exponent sum equals the full exponent
  sum, \(\sum_i\lambda_i^{+}=\sum_i\lambda_i\).
\end{lemma}
\begin{proof}\leanok
  \uses{exponents}
  The two finite sums differ only in the summands with \(\lambda_i\le 0\); under the hypothesis
  those are \(\lambda_i=0\), which contribute nothing to either side.
\end{proof}

\begin{theorem}[SRB reverse inequality, volume case]\label{sumPosExp_le_ksEntropy_of_SRB}
  \lean{ErgodicTheory.sumPosExp_le_ksEntropy_of_SRB}\leanok
  \uses{integral_log_abs_det_le_ksEntropy, sumPosExp_eq_sumAllExp_of_nonneg, sumAllExp_eq_integral_log_abs_det, SRBProperty, exponents, IsInjectivityPartition}
  Let \(T\) be ergodic on \(\R^d\) with a nonsingular, log-integrable derivative cocycle,
  differentiable, preserving an SRB (volume-case) measure (\(\mu\ll\operatorname{volume}\),
  \Cref{SRBProperty}), equipped with a nonempty injectivity partition \(\xi\) and integrable
  \(\log\rho,\ \log\abs{\det DT}\). If the Lyapunov spectrum is \emph{nonnegative}
  (\(0\le\lambda_i\) for all \(i\), the hypothesis \texttt{hspec}), then
  \[
    \sum_{\lambda_i>0}\lambda_i \;\le\; h_\mu(T) ,
  \]
  the reverse of the sharp Margulis--Ruelle bound (Pesin 1977; Ma\~n\'e 1981; Ledrappier--Young
  1985), discharged here via Rokhlin's inequality.
\end{theorem}
\begin{proof}\leanok
  \uses{integral_log_abs_det_le_ksEntropy, sumPosExp_eq_sumAllExp_of_nonneg, sumAllExp_eq_integral_log_abs_det}
  The chain is \(\sum\lambda^{+}\overset{(1)}{=}\sum\lambda\overset{(2)}{=}\int\log\abs{\det
  D_xT}\,d\mu\overset{(3)}{\le}h_\mu(T)\). Step (1) is
  \Cref{sumPosExp_eq_sumAllExp_of_nonneg}: the nonnegative spectrum kills the negative part. Step
  (2) is the trace--determinant identity (\Cref{sumAllExp_eq_integral_log_abs_det}) aligned to the
  Fr\'echet-derivative determinant. Step (3) is Rokhlin's inequality
  (\Cref{integral_log_abs_det_le_ksEntropy}).
\end{proof}

The spectrum gate \texttt{hspec} is \emph{exactly} the volume-case boundary and is not derivable
from \(\mu\ll\operatorname{volume}\). In the volume regime there are no genuinely negative
exponents, so \(\int\log\abs{\det DT}\,d\mu=\sum\lambda^{+}\) and the chain closes. For a
mixed-spectrum SRB measure with real stable directions,
\(\int\log\abs{\det DT}\,d\mu=\sum\lambda^{+}+\sum\lambda^{-}<\sum\lambda^{+}\), so Rokhlin's
inequality yields only the \emph{weaker} \(\sum\lambda^{+}+\sum\lambda^{-}\le h_\mu(T)\); recovering
the full reverse inequality then requires the Ledrappier--Young unstable-Jacobian machinery --- the
documented Mathlib-scale wall. (Volume-preserving hyperbolic maps have
\(\int\log\abs{\det}=0\) with genuinely negative exponents, so \texttt{hspec} is a real, separate
hypothesis, not a consequence of absolute continuity.)

\section{Pesin's entropy formula}

Combining the sharp Margulis--Ruelle upper bound with the volume-case reverse leaf by
antisymmetry gives Pesin's formula in the volume case, in both a clean spectral form and the
genuine integral form.

\begin{theorem}[Pesin's entropy formula, spectral form]\label{pesin_entropy_formula_spectral}
  \lean{ErgodicTheory.pesin_entropy_formula_spectral}\leanok
  \uses{margulisRuelle_sharp, sumPosExp_le_ksEntropy_of_SRB, SRBProperty, ksEntropy, exponents}
  For an ergodic differentiable self-map \(T\) of \(\R^d\) preserving an SRB (volume-case) measure
  \(\mu\) (\Cref{SRBProperty}) with nonsingular log-integrable derivative cocycle and nonnegative
  Lyapunov spectrum (\texttt{hspec}), and under the honest non-compactness atom-count input
  \texttt{hgeo} of the sharp Margulis--Ruelle inequality, the Kolmogorov--Sinai entropy equals the
  sum of the strictly positive Lyapunov exponents:
  \[
    h_\mu(T) \;=\; \sum_{\lambda_i>0}\lambda_i .
  \]
  This is Pesin's entropy formula in the absolutely continuous case (Pesin 1977;
  Ledrappier--Young 1985), obtained here foliation-free.
\end{theorem}
\begin{proof}\leanok
  \uses{margulisRuelle_sharp, sumPosExp_le_ksEntropy_of_SRB}
  Antisymmetry of the two inequalities. The \(\le\) direction
  \(h_\mu(T)\le\sum\lambda^{+}\) is the sharp Margulis--Ruelle inequality
  (\Cref{margulisRuelle_sharp}), which holds for \emph{every} invariant measure --- no SRB
  hypothesis --- and carries only its atom-count input \texttt{hgeo}. The \(\ge\) direction
  \(\sum\lambda^{+}\le h_\mu(T)\) is the volume-case reverse leaf
  (\Cref{sumPosExp_le_ksEntropy_of_SRB}). The hypotheses of the formula are exactly the union of
  the two halves.
\end{proof}

\begin{theorem}[Pesin's entropy formula, integral form]\label{pesin_entropy_formula}
  \lean{ErgodicTheory.pesin_entropy_formula}\leanok
  \uses{pesin_entropy_formula_spectral, UnstableJacobianRate}
  Under the same hypotheses, with \(\chi\) an unstable-Jacobian integrand
  (\Cref{UnstableJacobianRate}), the system entropy is the literal Pesin integral
  \[
    h_\mu(T) \;=\; \int_E \Bigl(\sum_{\lambda_i>0}\lambda_i\Bigr)\,d\mu
      \;=\; \int_E \chi\,d\mu .
  \]
\end{theorem}
\begin{proof}\leanok
  \uses{pesin_entropy_formula_spectral, UnstableJacobianRate}
  The bridge \(\int\chi\,d\mu=\operatorname{sumPosExp}\) holds because \(\chi\) is \(\mu\)-a.e.\
  equal to the constant \(\operatorname{sumPosExp}\) (\Cref{UnstableJacobianRate}) and \(\mu\) is a
  probability measure, so \(\int\chi\,d\mu=\operatorname{sumPosExp}\cdot\mu(E)=\operatorname{sumPosExp}\)
  via \texttt{integral\_congr\_ae}; no integrability hypothesis on \(\chi\) is needed. Rewriting
  through this bridge reduces the claim to the spectral form
  (\Cref{pesin_entropy_formula_spectral}).
\end{proof}

\paragraph{Vacuity disclosure (honest).} Exactly as for the expanding-map assembly
\Cref{pesin_formula_expanding}, these EuclideanSpace theorems are correct \emph{implications} whose
\emph{joint} hypothesis bundle has no known model on the non-compact space \(\R^d\). The bundle asks
for an ergodic, absolutely continuous \emph{probability} measure preserved by an
everywhere-nonsingular map with \emph{no} negative exponents (\texttt{hspec}); on \(\R^d\) these
clash --- an all-nonnegative spectrum expands volume on average, which on a non-compact space forces
mass to escape to infinity, incompatible with an a.c.\ invariant probability. So on \(\R^d\) the
statements are (as far as is known) vacuously-true assemblies of the implication, not exhibited
instances. The genuinely \emph{witnessed} equality lives on the compact circle, developed next.

\section{A witnessed instance: the doubling map}

The vacuity of the \(\R^d\) bundle is escaped on a compact carrier. For the doubling map
\(T:y\mapsto 2\cdot y\) on the unit circle \(\mathbb T=\operatorname{UnitAddCircle}\)
(\lean{ErgodicTheory.doublingMap}), ergodic for Haar measure
(\lean{ErgodicTheory.ergodic_doublingMap}), the binary partition
\(\alpha=\{[0,\tfrac12),[\tfrac12,1)\}\) (\lean{ErgodicTheory.Examples.Rokhlin.binPartition}) is a
genuine one-sided generator, so both sides of Pesin's formula are \emph{computed} and equal
\(\log 2\). This is the first non-vacuous full-system instance of the volume-case formula in the
library --- positive entropy on a compact space, every input honest.

\begin{theorem}[Per-partition Pesin identity for the doubling map]\label{pesin_identity_doublingMap_perPartition}
  \lean{ErgodicTheory.Examples.Rokhlin.pesin_identity_doublingMap_perPartition}\leanok
  \uses{ksEntropyPartition_doublingMap_eq_log_two, doublingMap_sumPosExp_eq_log_two}
  For the doubling map and its binary partition, the partition-relative Kolmogorov--Sinai entropy
  equals the sum of the positive Lyapunov exponents, both being \(\log 2\):
  \(h_\mu(T,\alpha)=\sum\lambda^{+}\).
\end{theorem}
\begin{proof}\leanok
  \uses{ksEntropyPartition_doublingMap_eq_log_two, doublingMap_sumPosExp_eq_log_two}
  The entropy side is \(h_\mu(T,\alpha)=\log 2\)
  (\Cref{ksEntropyPartition_doublingMap_eq_log_two}), the exponent side is
  \(\sum\lambda^{+}=\log 2\) (\Cref{doublingMap_sumPosExp_eq_log_two}); the two agree.
\end{proof}

\subsection*{The generator crux: dyadic arcs generate the Borel structure}

Upgrading the per-partition identity to the full system \(h_\mu(T)=h_\mu(T,\alpha)\) requires the
binary partition to \emph{generate}: its forward \(T\)-translates must recover every Borel set. This
reduces to showing the countable family of \emph{dyadic arcs} generates the Borel
\(\sigma\)-algebra, proved by a binary-expansion measurability argument.

\begin{definition}[The dyadic-arc family]\label{dyadicArcSet}
  \lean{ErgodicTheory.Examples.Rokhlin.dyadicArcSet}\leanok
  The set of all \(n\)-fold-join cells of the binary partition under the doubling map,
  \(\bigcap_{k<n}T^{-k}\bigl(\operatorname{binCell}(f(k))\bigr)\) over all lengths \(n\) and digit
  strings \(f:\Fin n\to\Fin 2\). Each such set is a dyadic arc of length \(2^{-n}\).
\end{definition}

\begin{definition}[Binary digit and dyadic partial sum]\label{binDigit}
  \lean{ErgodicTheory.Examples.Rokhlin.binDigit}\leanok
  The \(n\)-th binary digit \(d_n(y)=\mathbf 1_{T^{-n}(\operatorname{binCell}1)}(y)\in\{0,1\}\)
  records whether \(T^n y\) lands in the right half; the \(N\)-th dyadic partial sum
  \[
    \operatorname{binPartialSum}N(y)
      \;=\; \sum_{n=0}^{N-1} d_n(y)\,2^{-(n+1)}
  \]
  (\lean{ErgodicTheory.Examples.Rokhlin.binPartialSum}) is the truncated binary expansion of the
  canonical representative \(\operatorname{rep}(y)\in[0,1)\) (\lean{ErgodicTheory.Examples.Rokhlin.rep}).
\end{definition}

\begin{lemma}[One-step doubling recursion]\label{rep_doublingMap}
  \lean{ErgodicTheory.Examples.Rokhlin.rep_doublingMap}\leanok
  \uses{binDigit}
  On the representative, the doubling map acts by \(x\mapsto 2x\) on the left half and
  \(x\mapsto 2x-1\) on the right, uniformly
  \(\operatorname{rep}(Ty)=2\operatorname{rep}(y)-d_0(y)\).
\end{lemma}
\begin{proof}\leanok
  \uses{binDigit}
  Case split on \(y\in\operatorname{binCell}1\). On the right half \(d_0=1\) and
  \(2\operatorname{rep}(y)-1\in[0,1)\), which is the representative of \(2y\) after subtracting the
  integer \(1\) killed by the projection \(\R\to\mathbb T\); on the left half \(d_0=0\) and
  \(2\operatorname{rep}(y)\in[0,1)\) is already the representative.
\end{proof}

\begin{lemma}[Dynamical partial-sum recursion]\label{rep_eq_binPartialSum_add}
  \lean{ErgodicTheory.Examples.Rokhlin.rep_eq_binPartialSum_add}\leanok
  \uses{rep_doublingMap, binDigit}
  The representative splits into its first \(N\) digits plus a rescaled tail:
  \(\operatorname{rep}(y)=\operatorname{binPartialSum}N(y)+\operatorname{rep}(T^N y)\,2^{-N}\).
\end{lemma}
\begin{proof}\leanok
  \uses{rep_doublingMap}
  Induction on \(N\), feeding the one-step recursion \Cref{rep_doublingMap} into the tail
  (using that the \(0\)-th digit at \(T^N y\) is the \(N\)-th digit at \(y\)).
\end{proof}

\begin{lemma}[Convergence of the binary expansion]\label{tendsto_binPartialSum}
  \lean{ErgodicTheory.Examples.Rokhlin.tendsto_binPartialSum}\leanok
  \uses{rep_eq_binPartialSum_add}
  \(\operatorname{binPartialSum}N(y)\to\operatorname{rep}(y)\) as \(N\to\infty\).
\end{lemma}
\begin{proof}\leanok
  \uses{rep_eq_binPartialSum_add}
  The tail remainder \(\operatorname{rep}(T^N y)\,2^{-N}\le 2^{-N}\to 0\) is squeezed to zero, so
  the partial sums converge to \(\operatorname{rep}(y)\) by \Cref{rep_eq_binPartialSum_add}.
\end{proof}

\begin{lemma}[Digit sets are dyadic-measurable]\label{preimage_iterate_binCell_one_mem}
  \lean{ErgodicTheory.Examples.Rokhlin.preimage_iterate_binCell_one_mem}\leanok
  \uses{dyadicArcSet, binDigit}
  Each digit set \(T^{-n}(\operatorname{binCell}1)\) is the finite union of the generation-\((n+1)\)
  dyadic arcs whose last digit is \(1\), hence \(\operatorname{generateFrom}(\operatorname{dyadicArcSet})\)-measurable.
\end{lemma}
\begin{proof}\leanok
  \uses{dyadicArcSet}
  A point lies in \(T^{-n}(\operatorname{binCell}1)\) iff its length-\((n+1)\) digit string has last
  digit \(1\); realizing the string as the join cell exhibits the digit set as a finite union of
  arcs from \Cref{dyadicArcSet}.
\end{proof}

\begin{lemma}[The representative is dyadic-measurable]\label{measurable_rep_dyadic}
  \lean{ErgodicTheory.Examples.Rokhlin.measurable_rep_dyadic}\leanok
  \uses{binDigit, tendsto_binPartialSum, preimage_iterate_binCell_one_mem}
  \(\operatorname{rep}\) is \(\operatorname{generateFrom}(\operatorname{dyadicArcSet})\)-measurable.
\end{lemma}
\begin{proof}\leanok
  \uses{tendsto_binPartialSum, preimage_iterate_binCell_one_mem}
  Each digit function is the indicator of a dyadic-measurable digit set
  (\Cref{preimage_iterate_binCell_one_mem}), so each finite partial sum is dyadic-measurable; the
  pointwise limit \(\operatorname{rep}\) (\Cref{tendsto_binPartialSum}) inherits measurability by
  \texttt{measurable\_of\_tendsto\_metrizable'}.
\end{proof}

\begin{theorem}[Dyadic arcs generate the Borel structure]\label{borel_le_generateFrom_dyadicArcs}
  \lean{ErgodicTheory.Examples.Rokhlin.borel_le_generateFrom_dyadicArcs}\leanok
  \uses{dyadicArcSet, measurable_rep_dyadic}
  The Borel \(\sigma\)-algebra of the circle lies below the one generated by the dyadic arcs,
  \(m_{\mathbb T}\le\operatorname{generateFrom}(\operatorname{dyadicArcSet})\).
\end{theorem}
\begin{proof}\leanok
  \uses{measurable_rep_dyadic}
  Since \((\uparrow)\circ\operatorname{rep}=\operatorname{id}\) factors the identity through the
  measurable covering projection \(\R\to\mathbb T\), the Borel structure pulls back through
  \(\operatorname{rep}\) into \(\operatorname{generateFrom}(\operatorname{dyadicArcSet})\) by the
  dyadic-measurability of \(\operatorname{rep}\) (\Cref{measurable_rep_dyadic}); \(\operatorname{rep}\)
  is thus a Borel embedding and the containment follows.
\end{proof}

\begin{theorem}[The binary partition generates]\label{binPartition_isGenerating}
  \lean{ErgodicTheory.Examples.Rokhlin.binPartition_isGenerating}\leanok
  \uses{borel_le_generateFrom_dyadicArcs, dyadicArcSet, IsGenerating}
  The binary partition is a one-sided generator for the doubling map:
  \(\bigvee_{n}\operatorname{comap}(T^n)\,\sigma(\alpha)=m_{\mathbb T}\).
\end{theorem}
\begin{proof}\leanok
  \uses{borel_le_generateFrom_dyadicArcs, dyadicArcSet}
  The easy inclusion \(\le\) is immediate from measurability of the cells under the iterates. For
  \(\ge\): each digit set \(T^{-n}(\operatorname{binCell}i)\) is measurable for the saturated
  \(\sigma\)-algebra, hence each dyadic arc, being a finite intersection of digit sets, is too;
  reducing Borel to the dyadic-arc family by \Cref{borel_le_generateFrom_dyadicArcs} closes the
  reverse containment.
\end{proof}

\begin{theorem}[Entropy of the doubling map is \(\log 2\)]\label{ksEntropy_doublingMap_eq_log_two}
  \lean{ErgodicTheory.Examples.Rokhlin.ksEntropy_doublingMap_eq_log_two}\leanok
  \uses{binPartition_isGenerating, ksEntropy_eq_ksEntropyPartition_of_generating, ksEntropyPartition_doublingMap_eq_log_two}
  The Kolmogorov--Sinai entropy of the doubling map is \(h_\mu(T)=\log 2\).
\end{theorem}
\begin{proof}\leanok
  \uses{binPartition_isGenerating, ksEntropy_eq_ksEntropyPartition_of_generating, ksEntropyPartition_doublingMap_eq_log_two}
  The generator theorem (\Cref{ksEntropy_eq_ksEntropyPartition_of_generating}) collapses the
  system entropy to the binary partition, since it generates
  (\Cref{binPartition_isGenerating}); the per-partition entropy is \(\log 2\)
  (\Cref{ksEntropyPartition_doublingMap_eq_log_two}).
\end{proof}

\begin{theorem}[Doubling map: Pesin's formula, witnessed]\label{pesin_formula_doublingMap}
  \lean{ErgodicTheory.Examples.Rokhlin.pesin_formula_doublingMap}\leanok
  \uses{ksEntropy_doublingMap_eq_log_two, doublingMap_sumPosExp_eq_log_two, pesin_entropy_formula_spectral}
  For the doubling map,
  \[
    h_\mu(T) \;=\; \sum\lambda^{+} \;=\; \log 2 ,
  \]
  the genuine full-system Pesin equality on a compact carrier with strictly positive entropy. This
  is the first non-vacuous full-system instance of the volume-case Pesin formula
  (\Cref{pesin_entropy_formula_spectral}) in the library: the binary-expansion coding gives a
  genuine generating partition on a compact space, so both sides are computed rather than
  degenerate. Its companion \Cref{rokhlin_equality_doublingMap} exhibits the same equality against
  the integrated log-Jacobian \(\int\log\abs{\det DT}\,d\mu=\log 2\).
\end{theorem}
\begin{proof}\leanok
  \uses{ksEntropy_doublingMap_eq_log_two, doublingMap_sumPosExp_eq_log_two}
  The system entropy is \(\log 2\) (\Cref{ksEntropy_doublingMap_eq_log_two}) and the
  positive-exponent sum of the constant derivative cocycle \((2)\) is \(\log 2\)
  (\Cref{doublingMap_sumPosExp_eq_log_two}); the two agree.
\end{proof}
